\section{The Canonical Zeckendorf Fold and the Golden-Mean Language}
\label{sec:fold-language}

\subsection{Rotations, Windows, and the Folded Observable}

Let $\alpha\in(0,1)$ be irrational, let $\beta\in(0,1)$, and let
$x_0\in\TT:=\mathbb{R}/\mathbb{Z}$.  Define the Kronecker orbit
\[
x_t:=x_0+t\alpha \pmod 1,\qquad t\ge 0,
\]
and the binary readout
\[
s_t:=\mathbf{1}_{[0,\beta)}(x_t)\in\{0,1\}.
\]
For a fixed resolution $m\ge 1$, let
\[
\Omega_m:=\{0,1\}^m,\qquad
X_m:=\{u=(u_1,\dots,u_m)\in\{0,1\}^m:\ u_i u_{i+1}=0\text{ for }1\le i<m\}.
\]

\begin{definition}[Canonical fold]\label{def:canonical-fold}
Let \(F_1=F_2=1\) and \(F_{k+2}=F_{k+1}+F_k\). For
\(\omega=\omega_1\cdots\omega_m\in\Omega_m\), define its Fibonacci value by
\[
N(\omega):=\sum_{k=1}^m \omega_k F_{k+1}.
\]
By Zeckendorf's theorem \cite{Zeckendorf1972}, there is a unique sequence
\(Z(N(\omega))=(c_1,c_2,\dots)\in\{0,1\}^{\mathbb{N}}\) such that
\[
N(\omega)=\sum_{k\ge 1} c_k F_{k+1},
\qquad
c_kc_{k+1}=0\ \text{ for all }k\ge 1.
\]
Let
\[
\pref_m(c_1,c_2,\dots):=(c_1,\dots,c_m)\in X_m.
\]
The canonical fold at resolution \(m\) is
\[
\Fold_m(\omega):=\pref_m(Z(N(\omega)))\in X_m.
\]
\end{definition}

\begin{proposition}[Basic properties of the canonical fold]\label{prop:canonical-fold-basic}
For every \(m\ge 1\), the map \(\Fold_m:\Omega_m\to X_m\) is well defined,
surjective, and satisfies \(\Fold_m|_{X_m}=\mathrm{id}\). Moreover,
\(\Fold_m\) admits an order-independent local implementation by
value-preserving Zeckendorf rewrites.
\end{proposition}

\begin{proof}
Well-definedness is immediate from Definition~\ref{def:canonical-fold}. If
\(\omega\in X_m\), then \(\omega\) is already a legal Zeckendorf prefix, so the
uniqueness of the Zeckendorf expansion gives \(\Fold_m(\omega)=\omega\). Hence
\(\Fold_m\) fixes \(X_m\) pointwise, and surjectivity follows by taking
\(\omega=x\in X_m\) for any \(x\in X_m\). The order-independent local
implementation is recorded in Appendix~\ref{app:canonical-fold}.
\end{proof}

\begin{proposition}[Residue characterization of the canonical fold]\label{prop:fold-residue-characterization}
For every \(m\ge 1\) and \(\omega\in\Omega_m\), write
\[
N(\omega)=qF_{m+2}+r,
\qquad q\in\{0,1\},
\qquad 0\le r<F_{m+2}.
\]
Then
\[
\Fold_m(\omega)=\pref_m(Z(r)).
\]
Equivalently, for \(x\in X_m\),
\[
\Fold_m(\omega)=x
\iff
N(\omega)\equiv N(x)\pmod{F_{m+2}}.
\]
\end{proposition}

\begin{proof}
Since
\[
N(\omega)\le \sum_{k=1}^m F_{k+1}=F_{m+3}-2<F_{m+2}+F_{m+2},
\]
the quotient \(q\) is either \(0\) or \(1\). If \(q=0\), then \(r=N(\omega)\)
and the claim is immediate.

Assume \(q=1\). Then
\[
0\le r=N(\omega)-F_{m+2}\le F_{m+1}-2<F_{m+1},
\]
so the Zeckendorf expansion \(Z(r)\) has no nonzero digit in position \(m\) or
beyond. Hence appending a \(1\) in position \(m+1\) produces another legal
Zeckendorf digit string, namely the expansion of \(N(\omega)=F_{m+2}+r\). Its
first \(m\) digits are therefore exactly those of \(Z(r)\), which proves the
first statement.

If \(x\in X_m\), then
\[
N(x)\le \sum_{j=0}^{\lfloor (m-1)/2\rfloor}F_{m-2j+1}<F_{m+2},
\]
so \(x=\pref_m(Z(N(x)))\). Applying the first statement with
\(r\equiv N(\omega)\pmod{F_{m+2}}\) yields
\[
\Fold_m(\omega)=x
\iff
\pref_m(Z(r))=\pref_m(Z(N(x)))
\iff
r=N(x),
\]
which is equivalent to the congruence condition.
\end{proof}

\begin{definition}[Microstates and folded histograms]
For $t\ge 0$, the length-$m$ microstate is
\[
\omega_m(t):=s_t s_{t+1}\cdots s_{t+m-1}\in\Omega_m.
\]
The folded output is
\[
y_m(t):=\Fold_m(\omega_m(t))\in X_m.
\]
For a sample length $N$, define the empirical microstate law and folded
histogram by
\[
\hat\mu_{m,N}(u):=\frac{1}{N}\sum_{t=0}^{N-1}\mathbf{1}_{\{\omega_m(t)=u\}},
\qquad u\in\Omega_m,
\]
\[
\hat\pi_{m,N}(v):=\frac{1}{N}\sum_{t=0}^{N-1}\mathbf{1}_{\{y_m(t)=v\}},
\qquad v\in X_m.
\]
\end{definition}

Because irrational rotations are uniquely ergodic, the empirical microstate
law converges to a deterministic target law induced by Lebesgue measure.

\begin{definition}[Target laws]
Let $\mu_m$ be the length-$m$ microstate distribution induced by Lebesgue
measure on $\TT$:
\[
\mu_m(u):=\leb\{x\in\TT:\omega_m(x;0)=u\},\qquad u\in\Omega_m.
\]
Its folded pushforward is
\[
\pi_m^{\mathrm{fold}}:=(\Fold_m)_*\mu_m\in\Delta(X_m).
\]
\end{definition}

\begin{remark}[Spaces, maps, and laws]\label{rem:model-spaces-maps}
The basic objects may be summarized by
\[
\begin{array}{ccc}
\Omega_m & \xrightarrow{\Fold_m} & X_m\\[0.3em]
\mu_m \downarrow\phantom{(\Fold_m)_*} &  & \phantom{(\Fold_m)_*}\downarrow \pi_m^{\mathrm{fold}}\\[0.3em]
\hat\mu_{m,N} & \xrightarrow{(\Fold_m)_*} & \hat\pi_{m,N}
\end{array}
\qquad
\begin{array}{ccc}
\Omega_{m+1} & \xrightarrow{\Fold_{m+1}} & X_{m+1}\\[0.3em]
\pref_m \downarrow &  & \downarrow \rho_{m+1\to m}\\[0.3em]
\Omega_m & \xrightarrow{\Fold_m} & X_m.
\end{array}
\]
The left diagram separates microstate space \(\Omega_m\) from folded space
\(X_m\). The right diagram records the truncation interface whose failure to
commute is measured by the projective residual in
Appendix~\ref{app:discrepancy}.
\end{remark}

\begin{remark}[Scope]\label{rem:scope-fold-assumptions}
We fix the canonical fold throughout. The Sturmian rigidity and higher-block
conjugacy of Section~\ref{sec:sturmian-rigidity} use the Sturmian slice
$\beta=\alpha$. The sharp Parry divergence identity of
Section~\ref{sec:parry-divergence} is a statement about the Parry reference
family on $X_m$ and is independent of the fold once a law on $X_m$ has been
specified. The injective placement results of
Section~\ref{sec:injective-placement} combine the two: the low-complexity
support created by the folded rotation coding meets the ambient equilibrium
structure of the golden-mean shift.
\end{remark}

\subsection{The Parry Measure}

As a reference measure on \(X_m\) we use the Parry measure of the golden-mean
shift. Let
\[
A=\begin{pmatrix}1&1\\1&0\end{pmatrix},\qquad
\varphi=\frac{1+\sqrt5}{2}.
\]
The associated two-state Parry chain has transition matrix
\[
P=\begin{pmatrix}\varphi^{-1} & \varphi^{-2}\\ 1&0\end{pmatrix},
\]
and stationary distribution
\[
\pi(1)=\frac{1}{\varphi^2+1},\qquad
\pi(0)=1-\pi(1).
\]
For $u=u_1\cdots u_m\in X_m$, define the length-$m$ Parry cylinder law
\[
\pi_m(u):=\pi(u_1)\prod_{j=1}^{m-1}P_{u_j u_{j+1}}.
\]
We emphasize that \(\pi_m\) is used as a reference measure only. We do not
assume that \(\pi_m^{\mathrm{fold}}=\pi_m\) for irrational rotations.
Since \(|X_m|=F_{m+2}\), the ambient shift has topological entropy
\[
h_{\mathrm{top}}(X)=\lim_{m\to\infty}\frac1m\log |X_m|=\log\varphi,
\]
and the Parry measure is its unique measure of maximal entropy
\cite{LindMarcus1995,WaltersErgodicTheory}.

\begin{proposition}[Endpoint collapse of Parry cylinders]\label{prop:parry-endpoint-collapse}
For every legal word \(u=u_1\cdots u_m\in X_m\),
\[
\pi_m(u)=
\begin{cases}
\pi(0)\varphi^{-(m-1)}, & (u_1,u_m)=(0,0),\\
\pi(0)\varphi^{-m}, & (u_1,u_m)\in\{(0,1),(1,0)\},\\
\pi(0)\varphi^{-(m+1)}, & (u_1,u_m)=(1,1).
\end{cases}
\]
Equivalently,
\[
\pi_m(u)
=
\pi(0)\varphi^{-m+\mathbf{1}_{\{u_1=0\}}-\mathbf{1}_{\{u_m=1\}}}.
\]
In particular, the Parry cylinder weight depends only on the endpoints of the
word.
\end{proposition}

\begin{proof}
Let \(r_0:=\varphi\) and \(r_1:=1\). For every legal transition
\((i,j)\in\{(0,0),(0,1),(1,0)\}\), one has
\[
P_{ij}=\frac{r_j}{\varphi r_i}.
\]
Hence for a legal word \(u=u_1\cdots u_m\),
\[
\prod_{j=1}^{m-1}P_{u_j u_{j+1}}
=
\varphi^{-(m-1)}\frac{r_{u_m}}{r_{u_1}}.
\]
Multiplying by \(\pi(u_1)\) and using \(\pi(1)=\pi(0)\varphi^{-2}\) gives the
three cases above.
\end{proof}

\subsection{Sturmian Cylinder Bounds}

The following continued-fraction bounds on Sturmian cylinder lengths are used
in Sections~\ref{sec:sturmian-rigidity}
and~\ref{sec:injective-placement}.

\begin{theorem}[Bounds for Sturmian cylinder lengths]\label{thm:spg-sturmian-cylinder-measure-continued-fraction-sandwich}
Let $\alpha=[0;a_1,a_2,\dots]$ and let $p_n/q_n$ be its convergents.  If
$q_n\le t<q_{n+1}$, then every length-$t$ Sturmian cylinder $C(w)$ satisfies
\[
\frac{1}{q_{n+1}+q_n}
\;<\;
\leb(C(w))
\;<\;
\frac{1}{q_n}+\frac{1}{q_{n+1}}
\;<\;
\frac{2}{q_n}.
\]
Equivalently,
\[
\log\!\Bigl(\frac{q_n}{2}\Bigr)
\;<\;
-\log\leb(C(w))
\;<\;
\log(q_{n+1}+q_n).
\]
If $\alpha$ has bounded partial quotients, then
$-\log\leb(C(w))=\log t+O(1)$ uniformly in $w$.
\end{theorem}

\begin{proof}
The endpoints of length-$t$ cylinders are among the orbit points
$\{-k\alpha:0\le k\le t\}$, so cylinder lengths are the adjacent gaps in this
finite rotation set.  Standard continued-fraction gap estimates and the
three-gap theorem give
\[
\min_{1\le k\le t}\|k\alpha\|
=
\|q_n\alpha\|
\in
\left(\frac{1}{q_{n+1}+q_n},\frac{1}{q_{n+1}}\right),
\]
while the largest gap is at most
$\|q_{n-1}\alpha\|+\|q_n\alpha\|<q_n^{-1}+q_{n+1}^{-1}$; see
\cite{Slater1950DistributionIntegers,Slater1967GapsSteps}.  Taking negative
logarithms yields the second display.
\end{proof}
